\section{Tangent coupling with the full chiral defect projector}\label{sec:projectors}

\subsection{Field content and conventions}

The defect is $x^3=0$ in Euclidean spacetime $(x^0,x^1,x^2,x^3)$.
The bulk has a gauge field $A_\mu$, six adjoint scalar fields $X_A$,
$A=4,\ldots,9$, and their supersymmetric fermion partners. The triplets
$(X_4,X_5,X_6)$ and $(X_7,X_8,X_9)$ transform under $SU(2)_H$ and
$SU(2)_V$, respectively. Defect matter includes fundamental scalars
$q_m$, $m=1,2$, forming an H doublet, and fermions forming a V doublet;
$\bar q_m$ are the physical antifundamental conjugates. These are the
bulk/defect assignments of \cite{DFO,Baker}.

Work in a 32-dimensional complex Euclidean Clifford module, with
\begin{equation}\label{eq:chi-D}
 \chi=-\ii\Gamma_0\Gamma_1\cdots\Gamma_9,
 \qquad D=\Gamma_3\Gamma_4\Gamma_5\Gamma_6,
 \qquad \chi\epsilon=\epsilon,\quad D\epsilon=\epsilon.
\end{equation}
The second equation is the unfluxed D3--D5 intersection condition.
It is the projector appearing in \cite[eqs.~(2.7), (4.1)--(4.3),
(4.9)]{GW}; the terminating-brane boundary conditions and Nahm pole
are not part of our assumptions. The displayed conventions fix the
chirality and orientation signs. Complex ranks below are not counts
of real Euclidean Majorana--Weyl supercharges.

\subsection{A normal plane}

For the $(x^1,x^3)$-plane take the connection
\begin{equation}\label{eq:plane-connection}
 \mathcal B=\ii A_1\dd x^1+\ii A_3\dd x^3+X_8\dd x^1+X_6\dd x^3.
\end{equation}
This is the tangent prescription $n=(\dot x^1e_8+\dot x^3e_6)/|\dot x|$,
so the scalar direction changes with orientation. In the usual local
supersymmetry variation of scalar-coupled transport, a constant
Poincare parameter is annihilated by every tangent if
\begin{equation}\label{eq:plane-bps}
 [\dot x^1(\ii\Gamma_1+\Gamma_8)
  +\dot x^3(\ii\Gamma_3+\Gamma_6)]\epsilon=0.
\end{equation}
This is the tangent construction of \cite{Zarembo}, now with
\eqref{eq:chi-D} imposed as well. Define
\begin{equation}
 K_1=\ii\Gamma_8\Gamma_1,\qquad K_n=\ii\Gamma_6\Gamma_3.
\end{equation}

\begin{proposition}[Planar charge space]\label{prop:plane-charge}
The common subspace
$\chi=D=+1$, $K_1=K_n=-1$ has complex dimension two. Every parameter
in it obeys \eqref{eq:plane-bps} for every smooth planar tangent.
\end{proposition}
\begin{proof}
The four displayed matrices are commuting Hermitian involutions. Their
orthogonal projector is
\begin{equation}
 P_{\rm plane}=\frac1{16}(1+\chi)(1+D)(1-K_1)(1-K_n).
\end{equation}
Every nonidentity Clifford word in its expansion has zero trace, and
no nonempty product of these four generators is scalar. Hence
$\tr P_{\rm plane}=32/16=2$. Multiplication of $K_1\epsilon=-\epsilon$
by $\Gamma_8$, and of $K_n\epsilon=-\epsilon$ by $\Gamma_6$, gives
the two linear equations in \eqref{eq:plane-bps} separately.
\end{proof}

The triplet assignments are necessary within the stated real,
constant-map ansatz. The defect projector $D$ commutes with a defect
tangent gamma and with a V gamma; it anticommutes with the normal
gamma and with an H gamma. Pairing a tangent with H, or the normal
with V, therefore gives an involution anticommuting with $D$ and
no simultaneous eigenspinor. More generally, split a proposed real
scalar vector into H and V components. For an independently variable
defect tangent its wrong H component lies in the opposite
$D$-eigenspace from the remaining line equation, so it must annihilate
$\epsilon$ separately. Its Clifford square is its squared real length,
making it invertible unless that component is zero. The normal
argument exchanges H and V. Complex null internal vectors and
couplings restricted to one constrained tangent are outside this
argument. Proposition~\ref{prop:plane-charge} is the nonzero control.

\subsection{One charge across all defect directions}

To accommodate an angular family, introduce three tangent conditions
\begin{equation}
 K_\mu^{(s)}=\ii s_\mu\Gamma_{7+\mu}\Gamma_\mu,
 \qquad\mu=0,1,2,\quad s_\mu\in\{+1,-1\}.
\end{equation}
The corresponding oriented tangent isometry has determinant
$s_0s_1s_2$. Direct Clifford multiplication gives
\begin{equation}\label{eq:orientation-identity}
 K_0^{(s)}K_1^{(s)}K_2^{(s)}=(s_0s_1s_2)\chi D.
\end{equation}

\begin{proposition}[Orientation and the all-direction charge]\label{prop:all-charge}
For chirality $\chi=c$ and defect sign $D=d$, all three tangent
eigenvalues can be $-1$, together with $K_n=-1$, if and only if the
tangent isometry has determinant $-cd$. With that sign the common
space has complex dimension one; with the opposite sign it is zero.
\end{proposition}
\begin{proof}
Taking eigenvalues in \eqref{eq:orientation-identity} proves necessity.
For the allowed sign, one tangent condition follows from the other
two and $\chi,D$. The remaining five commuting conditions
$\chi=c$, $D=d$, $K_n=-1$, $K_0^{(s)}=K_1^{(s)}=-1$ are independent.
Their projector has trace $32/2^5=1$, by the same nonidentity-word
trace argument as above. Orthogonal changes of the oriented tangent
and internal frames give the general isometry statement.
\end{proof}

For $c=d=+1$, choose once and for all
\begin{equation}\label{eq:M-map}
 Me_0=-e_7,\quad Me_1=e_8,\quad Me_2=e_9,\quad Me_3=e_6.
\end{equation}
Then $M^TM=I$ and the connection is
\begin{equation}\label{eq:all-connection}
 \mathcal B=\ii A_\mu\dd x^\mu-X_7\dd x^0
              +X_8\dd x^1+X_9\dd x^2+X_6\dd x^3,
 \qquad\mu=0,1,2,3.
\end{equation}
The same nonzero constant complex parameter works for every Euclidean
four-dimensional tangent, hence for all normal planes indexed by
$S^2$. This is not a statement about arbitrary real Lorentzian
trajectories with real internal couplings.

\subsection{The remaining H weight}

Set $T_H=\ii\Gamma_4\Gamma_5$. The exact identity
\begin{equation}\label{eq:H-weight}
 D=K_nT_H
\end{equation}
implies $T_H\epsilon=-\epsilon$ for our signs. The tangent projectors
reduce the spacetime/V factors but do not introduce the other H weight.
Thus a surviving constant charge has the same one-dimensional H factor
at both endpoints. Using Baker's generator convention, the plus-normal
combination $\ii A_3+X_6$ is labeled by preserved generators $Q_{i2}$;
the conjugate convention exchanges 1 and 2. Spinor-parameter and
generator indices are related by charge conjugation, so the endpoint
labels must be fixed from the physical transformations, not read off
from a gamma eigenvalue alone.
